% !PRINTED 445
\weberpage{445}
durch $p$ theilbar sein, weil sowohl $\nu - \nu'$ als $\nu + \nu'$ kleiner als
$p$ ist. Unter den Zahlen

\weq{s}{1, 2, 3 \ldots p - 1}

kommt also die H\"alfte unter den Resten von $x^2$ vor, die H\"alfte
nicht. Die ersteren Zahlen und alle mit ihnen nach $p$ congruenten
heissen die \gesperrt{quadratischen Reste von $p$}, die anderen die
\gesperrt{quadratischen Nichtreste von $p$}.

Wenn nun $m$ quadratischer Rest ist, so ist die Congruenz
\[
x^2 \equiv m \ (\mathrm{mod}\ p)
\]
m\"oglich, und aus 7. und 8. folgt
\[
\left(\frac{m}{p}\right) = \left(\frac{x^2}{p}\right) = \left(\frac{x}{p}\right)^2 = +1.
\]
Es ist also $\left(\dfrac{m}{p}\right) = +1$, wenn $m$ quadratischer Rest von
$p$ ist.

Es ist aber in \S.\,135, 1. der Satz ausgesprochen, dass, wenn
f\"ur $m$ die Reihe der Zahlen 1, 2, \ldots p - 1 gesetzt wird,
$\left(\dfrac{m}{p}\right)$ ebenso oft das positive wie das negative
Zeichen hat, und daraus ersieht man, dass das negative Zeichen gilt,
wenn $m$ quadratischer Nichtrest ist, es folgt also:

11. \gesperrt{Ist $p$ eine ungerade Primzahl und $m$ durch $p$
nicht theilbar, so ist $\left(\dfrac{m}{p}\right) = +1$ oder $= -1$,
je nachdem $m$ quadratischer Rest oder quadratischer Nichtrest von
$p$ ist.}
Darin sind alle Haupts\"atze \"uber die quadratischen Reste
enthalten, z.\,B. der Satz:

Das Product zweier quadratischer Reste oder zweier Nichtreste
ist ein Rest; das Product aus einem Rest und einem Nichtrest ist
ein Nichtrest.

\begin{center}
\rule{0.25\linewidth}{0.4pt}
\end{center}
